\subsection{Thuật toán BYTE}

Thuật toán BYTE được đề xuất trong~\cite{zhang2022bytetrack} là một phương pháp
liên kết dữ liệu (data association) trong quá trình theo dõi đối tượng. BYTE
khác các thuật toán trước nó, vốn chỉ giữ lại các hộp phát hiện với độ tin cậy
cao, trao các hộp phát hiện độ tin cậy thấp cơ hội sau khi đã xem xét các hộp
độ tin cậy cao. Cách hoạt động như sau: đầu tiên, thuật toán nhận vào hộp phát
hiện được dự đoán trong khung hình hiện tại và chia nó thành hai phần bằng cách
lấy ngưỡng trên độ tin cậy; các hộp có độ tin cậy cao được đem liên kết với các
vết theo dõi được dự đoán theo Kalman Filter từ khung hình trước bằng độ đo
tương đồng 1 qua Thuật toán Hungarian; sau đó, các hộp có độ tin cậy thấp được
đem liên kết với các vết theo dõi còn lại theo độ đo tương đồng 2; cuối cùng,
các hộp có độ tin cậy cao chưa được liên kết sẽ được tạo thành các vết theo dõi
mới~\cite{zhang2022bytetrack}. Đây là một phương pháp liên kết dữ liệu đơn
giản, hiệu quả và rất tổng quát. Ta có thể dùng bất cứ độ đo nào trong hai lần
liên kết như IoU hay Re-ID.

Trong bài báo gốc, tác giả sử dụng chung một ngưỡng chọn các hộp độ tin cậy cao
và tạo ra các hộp theo dõi mới. Tuy nhiên, trong mã nguồn thực tế, hệ thống sẽ
chọn hai giá trị khác nhau cho hai ngưỡng này nhằm tăng tính linh hoạt và ổn
định.
